% !TEX root = ../main.tex
\vspace{-5mm}
\section{Preliminaries}
\vspace{-3mm}
\ecoparagraph{Notations.}
  We represent scalar discrete random variables over $N$ categories as one-hot column vectors, with ${\VC = \{x \in \{0,1\}^N \colon \sum_{i=1}^N x_i = 1\}}$ denoting the set of such vectors. We further define the probability simplex $\Delta \vcentcolon= \{ z \in \RR^N\colon z \ge 0, \langle z, \vec{1} \rangle = 1 \}$ as the convex hull of $\VC$, where $\vec{1} \in \RR^N$ denotes the all-ones vector. Additionally, let $I \in \RR^{N \times N}$ be the identity matrix. The matrix-weighted inner product is defined as $\langle a, b \rangle_Q \vcentcolon= a^\top Q b$, and we omit the subscript when $Q = I$. For absorbing diffusion process we define special mask token $m \in \VC$ such that $m_N = 1$.

\vspace{-2mm}
\subsection{Diffusion Language Models}
\label{sec:discrete diffusion models}
\vspace{-3mm}
  We start by recalling the Diffusion Language Models (DLMs; \citealp{austin2021structured, campbell2022continuous, lou2024discrete}). Consider a distribution $p^*(x_0) \in \Delta$ over $\VC$. DLMs aim to generate samples from $p^*(x_0)$ by constructing a transformation process from an initial tractable distribution, based on predefined \textit{forward diffusion}.

\ecoparagraph{Forward Diffusion.}
  A forward discrete diffusion process is defined by a family of time-dependent distributions $p_t(x_t) \in \Delta,$ evolving over the interval $t \in [0, 1].$ This evolution follows a continuous-time Markov chain defined by a linear ordinary differential equation~\citep{anderson2012continuous, campbell2022continuous, lou2024discrete}:
\vspace{-2mm}
  \begin{equation}
    \frac{dp_t}{dt} = Q_t p_t, \quad p_0 = p^*,
  \label{eq:forward diffusion}
    \vspace{-0.5mm}
  \end{equation}
  where $Q_t \in \RR^{N \times N}$ are the diffusion matrices and have nonnegative non-diagonal entries and columns which sum to zero, ensuring valid probability dynamics. These matrices are typically chosen in the form $Q_t = \sigma_t Q$, where $Q$ is a fixed base matrix that defines the structure of the diffusion process, and $\sigma_t$ is a time-dependent scalar that controls the diffusion schedule. With an appropriate choice of $Q$ and scheduling function $\sigma_t$, the distribution $p_t$ converges to a tractable limiting distribution $\pi(x_1) \in \Delta$ as $t \to 1$.

  For a fixed initial state $x_0 \sim p^*(x_0)$, the forward diffusion process defined in equation~\eqref{eq:forward diffusion} induces a family of conditional distributions over intermediate states $x_t \in \VC$.
%  \begin{Box1}{Conditional Distribution}
    \begin{equation}
      p_{t \mid 0}(x_t \mid x_0) \vcentcolon= \text{Cat}(x_t; \exp\left(\bar{\sigma}_t Q\right)x_0 ), 
    \label{eq:conditional_distribution}
    \end{equation}
%  \end{Box1}
  where, $\exp\left(\bar{\sigma}_t Q\right)$ denotes the matrix exponential,
and $\bar{\sigma}_t \coloneqq \int_{0}^{t} \sigma_s \, ds$ is the cumulative schedule.


\ecoparagraph{Concrete Score Matching (SEDD).}

  The goal of the Concrete Score Matching~\citep{meng2022concrete, lou2024discrete} is to approximate the ratio of
marginal probabilities $\frac{p_t(y)}{p_t(x)}$. To achieve this, \citep{lou2024discrete} propose learning a score
function $\widehat{s} \colon \mathcal{V} \times [0, 1] \to \mathbb{R}_{+}^N$ using the
\textit{Score Entropy Discrete Diffusion} (SEDD) loss. This objective, defined for a data point $x_0 \sim p^*(x_0)$, is:
  \begin{Box1}{SEDD Loss}
    \vspace{-8pt}
    \begin{gather}
    \label{eq:sedd loss}
       \!\!\!\!\LC_{\text{SEDD}}(\widehat{s}, x_0) \colon\!\!\!\!= \!\!\!\int_0^1\!\! \EE_{p_{t|0}(x_t|x_0)} \biggl[\sum_{y \neq x_t} \lambda_{x_ty}
    \!\!\times\!\! \biggl(\!\!\langle \widehat{s}(x_t,t), y\rangle - \frac{p_{t \mid 0}(y \mid x_0)}{p_{t \mid 0}(x_t \mid x_0)} \log  \langle \widehat{s}(x_t, t), y\rangle \biggr)\!\!\biggr]dt,
    \end{gather}
  \end{Box1}
  where $\lambda_{x_ty}$ is positive weighting function.
  \citet{lou2024discrete} show that the minimizer for the data distribution $p^*(x_0)$: 
  \[
    s^*= \arg\min\nolimits_{\widehat{s}} \EE_{p^*(x_0)}\left[\LC_{\text{SEDD}}(\widehat{s}, x_0)\right]
  \] 
  recovers the desired ratio of marginal probabilities $\frac{p_t(y)}{p_t(x)}.$ 
  
  However, as noted in~\citep{lou2024discrete}, applying this formulation to large-scale settings (e.g., GPT-2 with $N = 50257$ tokens) is practically challenging, as storing and computing all edge weights in the transition matrix $Q_t$ becomes infeasible. To address this, subsequent work has focused on specific classes of diffusion processes with simplified transition structures, reducing computational overhead while improving numerical stability and model efficiency.

\ecoparagraph{Absorbing Process (MDLM).}

  Previous work on DLMs has primarily focused on two particular classes of processes. The first is the \textit{absorbing process} (see Appendix~\ref{app:absorbing process} for further details). A prominent example of that is the \textit{Masked Diffusion Language Model} (MDLM; \citealp{sahoo2024simple, shi2024simplified}), in which the authors, rather than directly parameterizing the probability ratios $\frac{p_t(y)}{p_t(x)}$, follow the strategy of~\citet{ho2020denoising, austin2021structured, campbell2022continuous} and instead learn the reverse conditional density $p_{0\mid t}(x_0 \mid x_t)$. Accordingly, they adopt an alternative parameterization $\widehat{x}_0\colon \VC \times [0, 1] \to \Delta$, which yields the following training objective for a fixed data point $x_0~\sim~p^*(x_0)$:
  \begin{Box1}{MDLM Loss}
    \vspace{-8pt}
    \begin{gather}
      \textstyle \LC_{\text{MDLM}}(\widehat{x}_0, x_0) \vcentcolon= \int_0^1 \, \EE_{p_{t|0}(x_t \mid x_0)} \left[
          \lambda_{t} \langle \log\widehat{x}_0(x_t, t), x_0 \rangle
      \right]dt,
    \label{eq:mdlm loss}
    \end{gather}
  \end{Box1}
  %
  where $\lambda_{t}$ is positive weighting function and $\log \widehat{x}_0(x_t, t)$ is element-wise logarithm of $\widehat{x}_0(x_t, t)$. Note that in the original work, the authors use the notation $\log \langle \widehat{x}_0(x_t, t), x_0 \rangle$, which is equivalent. Moreover, \citep{sahoo2024simple} show that this loss corresponds to the negative evidence lower bound (NELBO). As a result, at optimality, the learned model on whole data distribution $p^*(x_0)$: 
  \[
    x_0^* = \arg\min_{\widehat{x}_0} \EE_{p^*(x_0)}\left[\LC_{\text{MDLM}}(\widehat{x}_0, x_0)\right]
  \]
  recovers the target reverse conditional density $p_{0\mid t}(x_0 \mid x_t).$

\ecoparagraph{Uniform Process (UDLM/Duo).}

  Another important case is the \textit{uniform process} (see Appendix~\ref{app:uniform process} for further details). Two notable models have been proposed in this setting: the \textit{Uniform Diffusion Language Model} (UDLM; \citealp{schiff2024discreteguidance}) and its extension, \textit{Diffusion Duality} (Duo; \citealp{sahoo2025the}). UDLM adopts the same model parameterization and target function as MDLM, and introduces the loss for a data point $x_0 \sim p^*(x_0)$ as
  \begin{equation}
    \mathcal{L}_{\text{UDLM}}(\widehat{x}_0,x_0) \vcentcolon= \textstyle \int_0^1 \EE_{p_{t|0}(x_t \mid x_0)} g\bigl(x_t, x_0, \widehat{x}_0(x_t, t)\bigr) dt,
  \label{eq:udlm loss}
  \end{equation}
  where the function $g\bigl(x_t, x_0, \widehat{x}_0(x_t, t)\bigr) \vcentcolon=$
  \begin{gather}
       \lambda_{t}\Bigg(\frac{1}{p_{t\mid0}(x_t \mid x_0)} - \frac{1}{p_{t \mid 0}(x_t \mid \widehat{x}_0)}
          - \sum_{y \neq x_t} \frac{p_{t \mid 0}(y \mid x_0)}{p_{t\mid0}(x_t \mid x_0)}\log\frac{p_{t \mid 0}(x_t \mid \widehat{x}_0) \, p_{t \mid 0}(y \mid x_0)}{p_{t \mid 0}(y \mid \widehat{x}_0) \, p_{t \mid 0}(x_t \mid x_0)}\Bigg), \nonumber
  \end{gather}
  with $\lambda_t$ denoting a positive weighting function. For clarity, we slightly abuse notation by omitting
the model's input in $\widehat{x}_0$, assuming $\widehat{x}_0 = \widehat{x}_0(x_t, t)$ within the expression.
  Duo further introduces an objective with reduced variance. To achieve it the authors first introduce distribution:
%  \begin{equation}
$
    \tilde{p}_{t\mid 0}(w_t \mid x_0) = \NC(\tilde{\alpha}_t x_0, (1-\tilde{\alpha}_t^2)I), \nonumber
$
%  \end{equation}
  where $\tilde{\alpha}_t$ denotes a rescaled diffusion schedule. Then they show that for $w_t \sim \tilde{p}_{t\mid 0}(w_t \mid x_0)$ and $x_t(w_t) \vcentcolon= \arg\max(w_t)$:
\vspace{-0.5mm}
\[
    x_t(w_t) \sim p_{t|0}(x_t \mid x_0),
  \vspace{-1mm}
  \]
  where $\arg\max(w_t)$ denotes a one-hot vector whose index corresponds to the maximum value in $w_t$. Based on this, they propose reducing the variance of the objective by substituting the discrete model input $x_t(w_t)$ with soft approximation:
  \begin{equation}
    x^{\tau}_{t}(w_t) = \mathrm{softmax}\left({w_t}/{\tau}\right). \nonumber
    \vspace{-2mm}
  \end{equation} 

  This relaxation leads to an extended model parameterization $\widehat{x}_0 \colon \Delta \times [0, 1] \to \Delta$. The resulting loss function for a data point $x_0 \sim p^*(x_0)$ is given by:
  \begin{Box1}{Duo Loss}
    \vspace{-8pt}
    \begin{gather}
    \label{eq:duo loss}
      \textstyle \LC_{\text{Duo}}(\widehat{x}_0, x_0) \vcentcolon= \int_0^1 \EE_{\tilde{p}_{t\mid 0}(w_t\mid x_0)} \left[g(x_t(w_t),x_0, \widehat{x}_0(x_t^{\tau}(w_t), t))\right]dt.
    \end{gather}
  \end{Box1}
  %
  Moreover, \citet{sahoo2025the} shows that, in the limit as $\tau \to 0^{+}$, this loss corresponds to the NELBO of UDLM. Consequently, at optimality and in this limit, the learned model on whole data distribution $p^*(x_0)$:
  \[
    x_0^* = \arg\min_{\widehat{x}_0} \EE_{p^*(x_0)}\left[\LC_{\text{Duo}}\right](\widehat{x}_0, x_0)
  \]
  recovers the reverse conditional distribution $p_{0\mid t}(x_0 \mid x_t).$
  
\ecoparagraph{Sequence-Level Discrete Diffusion.}
  In practical applications, we are interested in generating sequences of some length. A direct extension of the diffusion process to this setting would require modeling over an exponentially large space, resulting in a diffusion matrix of intractable size. To overcome this limitation, following~\citep{lou2024discrete}, we consider diffusion processes defined by sparse diffusion matrices that perturb tokens independently across sequence positions. Under this assumption, the forward diffusion matrix factorizes across tokens, yielding both a factorized conditional distribution and, crucially, a loss function that decomposes as a sum of token-wise losses. As a result, the per-token objectives in equations~\eqref{eq:sedd loss}, \eqref{eq:mdlm loss}, and~\eqref{eq:duo loss} can be directly applied to sequence modeling.

\vspace{-1mm}
\subsection{Distillation of Diffusion Models} 
\vspace{-2mm}
\ecoparagraph{Inverse objective.}
  We build on the \emph{Inverse Distillation} framework~\citep{gushchin2025inverse, kornilov2025universal}, which optimizes a \emph{student distribution} $p_\theta(x_0)$ parametrized by a few-step generator such that the fixed teacher $f^*$ remains optimal under $p_\theta$. Specifically, it considers the objective
  \begin{gather}
   \LC_{\text{inv}}(\theta) \vcentcolon= \!\!\! \EE_{p_{\theta}(x_0)} \!\left[\LC_{\text{cont.}}(f^*\!\!, x_0)\right]  - \min\nolimits_{\widehat{f}} \EE_{p_{\theta}(x_0)} \! \left[\LC_{\text{cont.}}(\widehat{f}, x_0)\right],
  \label{eq:inverse optimization problem}
    \vspace{-1mm}
  \end{gather}

  where the inner minimizer is the \emph{fake} model trained on samples from $p_\theta$. 
  By construction, the inverse loss satisfies $\LC_{\text{inv}}(\theta) \geq 0$ and for $p_\theta(x_0) = p^*(x_0)$ the minimum value $0$ is attained. However, this formulation does not, in general, guarantee the uniqueness of the minimizer meaning that optimization may converge to suboptimal solutions.

  While such a general description of the inverse distillation was proposed in~\citep{kornilov2025universal} the particular cases of this distillation were introduced before, specifically for continuous diffusion (Score identity Distillation, SiD) in~\citep{zhou2024score} and specifically for continuous flows (Flow Generator Matching, FGM) in~\citep{huang2024flow}. 

  \ecoparagraph{Connection to this work.}
    In \emph{continuous} space, the model distribution $p_\theta$ is typically represented via a few-step generator and reparameterized, allowing the optimization of~\eqref{eq:inverse optimization problem} via backpropagation through sampled trajectories. In contrast, extending this approach to the \emph{discrete} domain presents several challenges, stemming from the inherent characteristics of discrete spaces. These include \underline{theoretical concerns} regarding the validity of the training objective, due to the non-uniqueness of its minimizers, as well as \underline{practical difficulties} related to the non-differentiability of discrete sampling. Our objective is to extend the inverse distillation framework to the widely studied class of \emph{diffusion language models} by addressing both of these challenges, as discussed in \wasyparagraph\ref{sec:idlm}.
